\introduction % Структурный элемент: ВВЕДЕНИЕ

В современном мире актуальна проблема соблюдения конфиденциальности при совместной обработке данных. При совместной работе нескольких организаций задача состоит в получении общего результата без раскрытия приватной информации другим участникам.

Для совместной обработки приватных данных применяются многосторонние конфиденциальные вычисления (МКВ; англ. Multi-\allowbreak Party Computation). Классические модели МКВ позволяют двум или более участникам совместно вычислить значение функции на объединённых входных данных, не раскрывая сами данные и гарантируя корректность результата.

Сферы применения МКВ включает совместную аналитику \cite{17}, финансы и здравоохранение. Для задач МО необходимо сочетать защиту данных с требованиями к точности, производительности и масштабируемости вычислений.

Цель работы --- анализ теоретических и криптографических основ МКВ, классификации протоколов и их применений, а также планирование разработки продукта с применением МКВ в области МО.

Для достижения цели поставлены следующие задачи:

\begin{enumerate}

\item Анализ теоретических основ и свойств безопасности МКВ.

\item Исследование криптографических примитивов, используемых в МКВ.

\item Классификация протоколов МКВ.

\item Анализ практических применений МКВ.

\item Выявление ограничений МКВ, условий их практического использования, а также определение направлений развития МКВ.

\item Разработка плана создания и оценки прототипа защищенного МО.

\end{enumerate}

Теоретические основы изложены в разделе~\ref{sec:theoretical-foundations}, криптографические примитивы --- в разделе~\ref{sec:cryptographic-foundations}, классификация --- в разделе~\ref{sec:protocol-classification}, применения --- в разделе~\ref{sec:applications}, ограничения и перспективы --- в разделе~\ref{sec:limitations}, описание проекта защищенного МО --- в разделе~\ref{sec:future-research}; заключение содержит итоги и рекомендации.